% QMB_n15_zeit_superposition.tex — Chapter 15: Time in State Space, Superposition Without Time
% Substantive basis: Doc. 307 (July 2026), Doc. 334 (Aug. 2026), Doc. 333; own presentation
\chapter{Where Time Sits: State Space and Superposition}
\label{ch:zeit}

\section{The Bifurcation}

At the beginning of every quantum mechanical model building stands a choice that
usually remains unspoken and nevertheless determines everything else: Is
time a direction \emph{in} state space, or does it run
\emph{alongside} it, as an external parameter?

The question sounds technical, but is ontological. It does not decide
which mathematics is used --- both answers work with
self-adjoint operators, spectra, and projectors ---, but \emph{where}
time sits. And precisely this placement decides whether the
physical quantities \emph{fall} out of the structure or must be
\emph{attached} to it.

The question reaches back to the roots of quantum theory. As early as 1932,
Wolfgang Pauli recognized that a time operator canonically conjugate to the
Hamiltonian encounters fundamental difficulties. The Pauli theorem has
shaped the development of quantum mechanics and seemingly made time as an external parameter
irrevocable. Later works --- Aharonov and Bohm, Garrison and Wong ---
showed that under certain conditions, such as for non-
bounded-below Hamiltonians or compact configurations,
time operators can nevertheless be constructed. This exception is, for the FFGFT,
the rule.

\section{The Unrolled View: Time as External Parameter}

The familiar formulation situates time outside the state space.
The Hilbert space is the space of states \emph{at one
time}, typically $\mathcal{H} = L^2(\text{configuration space})$.
Time is the parameter $t$ along which the state
evolves, mediated by $e^{-i\hat Ht}$. Position, momentum, spin are
operators on $\mathcal{H}$. Time is none.

That it cannot be one is not an oversight, but a result:
Pauli's objection shows that for a Hamiltonian bounded from below,
no self-adjoint, canonically conjugate
time operator exists.\status{E} In the unrolled view, time therefore necessarily remains the stage, not the play.

The consequence is structural: Since time stands outside, the state space
alone contains no dynamical scales. Masses, couplings,
lifetimes are not present in the pure spatial structure --- they
must be introduced from outside. In the Standard Model, this happens
through Yukawa terms and the free parameters that fix their strength. The structure provides the framework; the numbers come from
measurement.

The problem is old. Weyl and Wigner recognized that the masses of
elementary particles appear group-theoretically as Casimir invariants
--- as structural labels, not as attached parameters.
Nevertheless, the origin of the mass values remained open. The Higgs mechanism explains
the mass of the gauge bosons, but introduces new free parameters itself.
The criticism of the unrolled view is not that it is wrong,
but that it leaves the origin of the mass scales unanswered and
refers to external inputs.

\subsection*{A Representative Construction of This Kind}

Many model programs that strive for a geometric origin of
particle properties follow, despite all differences, the same
scheme: One chooses a discrete carrier (a lattice, a
cut-and-project structure, a periodic graph); defines on it
a self-adjoint operator family; extracts sectors via
spectral projectors; and assigns physical meaning to these sectors
afterwards --- charge, spin, mass, the latter usually via
a parametric mapping with one or two free parameters.

Such a procedure can be mathematically fully controlled.
The decisive observation concerns not its rigor, but
its placement of time: The carrier is a spatial object, the
operator evolution runs over an external parameter. The sectors are
spatial eigenspaces, not time modes. Mass is not
contained in them; it is attached.

And here an argument applies that carries itself: \textbf{A
parametric mapping with free parameters can hit any finite
list of values.} If the mapping is not fixed in advance and universally,
it possesses no predictive content --- it renames the
known inputs. To deliver numbers, such a construction must
be coupled to external values: to measurement, and
thus indirectly to the Standard Model itself. That is honest, provided it declares its limit:
that it does not derive masses, but provides a
hypothesis architecture. The difficulty of this path is
not accidental, but the consequence of the unrolled placement of time.

A concrete example: Graph modeling defines
quasiparticles on hexagonal lattices. The excitations show linear
dispersion and can be read as massless fermions --- but the
introduction of a mass always requires additional terms (Haldane term,
Kekulé distortion) with free parameters. The geometry provides the
kinematics, not the mass scale.

\section{The Rolled-Up View: Time as Compact Dimension}

The alternative does not leave time outside and does not eliminate it either,
but pulls it \emph{into} the state space --- as a closed,
periodic dimension. In place of $L^2(\text{space})$ with external
$t$ comes a space in which a compact time direction is part of the
geometry: $\mathcal{H} = L^2(T^4)\otimes\mathbb{C}^3$, where one of the
torus directions carries time. The bijectivity of this representation
(Type III) is documented in the corpus (Chapter~\ref{ch:hilbert}).

On a closed dimension, the admissible states are standing
modes with discrete winding numbers. Quantization is then not
postulated, but geometrically enforced: Only certain modes fit
on the closed structure. Mass becomes an eigenvalue of an
operator on $\mathcal{H}$ itself --- not attached, but
structurally contained. The ratios of the masses are
fixed, because the admissible modes are fixed.\status{B}

The relation $\tilde T\cdot m = 1$ is the concrete execution of this
idea: Every massive quantity \emph{is} its own time cycle. The
decompactification $S^1_m\to\mathbb{R}_t$ is measure-preserving, thus
lossless. That a compact dimension generates a discrete mode tower
whose masses behave as the inverse of the compactification period
is not an exotic feature, but the mechanism that Kaluza
and Klein demonstrated for a compact spatial dimension --- here applied to
time.

Pauli's objection does not hit the rolled-up view. It applies to an
open time unbounded below conjugate to a bounded
Hamiltonian. A compact time behaves like an
angle variable: Its conjugate spectrum is discrete, as with
angular momentum to the angle, and the Pauli argument does not apply.\status{B} The
closed time is not only admissible, it is the natural place
of a discrete mode structure.

This insight has roots: Witten pointed out that compact
time dimensions occur in string theory; Hawking discussed the
implications of closed time for cosmology; Bars developed
two-time physics. And the rolled-up view touches
thermodynamic time: If time is compact, the arrow direction is not
fundamental, but an emergent property of the non-compact
spatial parts --- a natural explanation of macroscopic
irreversibility without assumptions about initial conditions.

\begin{laierspur}
One can build a clock in two ways. Either one hangs a
pendulum on a frame and lets time run from outside --- then
one must adjust the length of the pendulum by hand. Or one builds the
clock as a closed circle in which each oscillation returns into itself
--- then the circle itself determines which oscillations are possible.
The FFGFT builds the clock in the second way. Therefore it does not need
to adjust the masses: They fall out of the geometry.
\end{laierspur}

\section{The Core of the Trade-off}

\textbf{What the unrolled view gains:} the seamless connection to
the machinery of quantum field theory; the freedom from a strong
ontological commitment; the familiarity of the open time axis.
\textbf{What it costs:} The structure alone does not fix any masses.
Dynamical values must come from outside; predictions depend on
external calibration; a model that understands itself as independent
couples for its numbers back to measurement.

\textbf{What the rolled-up view gains:} The modes and their
ratios fall out of the structure, without parametric attachment; the
discrete quantization is a gift rather than a postulate. \textbf{What it
costs:} a strong ontological claim --- time as a closed
dimension --- that contradicts the accustomed intuition of running time.
And it does not eliminate every external connection: The absolute
scale still needs a unit bridge, the SI conversion via fixed
conversion factors.

But the decisive finding is this: \textbf{Both paths use
the same operator formalism.} The difference lies not in the
mathematics, but solely in whether time sits inside or outside.
One can keep the same construction --- self-adjoint operator,
spectral sectors --- and only pull time into the space;
then the sectors become modes, and the mass that previously had to be attached
becomes an eigenvalue of the structure. The choice of
time placement is a framework choice that precedes model building and
determines its interpretive reach.

\section{Who Has Already Pulled Time Inside}

The idea of not treating time as an external parameter is
not a novelty, but a long-followed line. It is honest to
locate the rolled-up view in this line and to delimit its peculiarity.

\textbf{Relational time.} Page and Wootters let time arise within
the state space: as entanglement between a clock subsystem and the rest;
the total state is stationary,
\enquote{evolution without evolution}. Moreva et al. and Marletto and
Vedral have refined and experimentally illustrated the mechanism;
Gambini and Pullin developed the Montevideo interpretation.

\textbf{Timeless quantum gravity.} Wheeler--DeWitt: The universal
wave function is timeless, $\hat H\Psi = 0$; the resulting
\enquote{problem of time} is the statement that there is no external
time. Kuchař and Isham provided the overviews. Rovelli:
relational quantum mechanics, with Connes the thermal time hypothesis.
Barbour: the consistent elimination of time in favor of timeless
configurations.

\textbf{Compact time in string theory.} Bars and colleagues with
two-time physics and extended symmetry group.

\textbf{The mode side.} Wigner: Mass is a Casimir invariant, a
label of irreducible representations of the Poincaré group --- already a
structural eigenvalue. Kaluza and Klein: A compactified
dimension generates a discrete mode tower with quantized charges
--- the blueprint for the fact that a rolled-up dimension brings forth modes.

The demarcation is precise: The majority of this tradition
makes time relational, emergent, or timeless --- it takes time
as fundamental \emph{out}. The rolled-up view takes time
\emph{in}: as a compact geometric dimension that carries modes.
It shares with Page--Wootters and Rovelli the rejection of the external
parameter time, but differs in the positive proposal.

\section{Mathematical Execution}

The Hilbert space of the rolled-up view is
$\mathcal{H} = L^2(T^1\times M)\otimes\mathcal{H}_{\mathrm{int}}$, with
a compact time dimension $T^1$ of length $T$ and the ordinary
space $M$. A function on it expands in a Fourier series:
\begin{equation}
  \Psi(t,\mathbf{x}) = \sum_{n\in\mathbb{Z}}e^{i n t/T}\,\psi_n(\mathbf{x}).
\end{equation}
The discrete frequencies $\omega_n = 2\pi n/T$ correspond to the
mass values in natural units. The operator that generates them
is $\hat E = -i\,\partial_t$ with the spectrum $\{2\pi n/T\}$. The
quantization of masses is a direct consequence of the periodicity
of time; nothing needs to be postulated.

\textbf{Yukawa coupling as selection rule.} The coupling between modes
is a matrix element of the interaction operator:
\begin{equation}
  g_{ijk}\propto\int_{T^1}e^{i(n_i-n_j-n_k)t/T}\,dt
  = T\cdot\delta_{n_i,\,n_j+n_k}.
\end{equation}
This is an exact conservation of the mode index in the time direction. The
Yukawa couplings of the Standard Model appear as selection rules for
the mode indices, not as free parameters.\status{B}

\textbf{Charge quantization.} The compact time direction behaves
like a fifth dimension in the sense of Kaluza--Klein, but timelike.
The charge is $Q\propto n/R_T$ with the radius $R_T$ of the compact time
--- charge quantization as a consequence of periodicity.

\textbf{Mass hierarchy.} The corpus describes the lepton masses in
three mutually translatable representations: first, fundamentally via
quantum numbers $(n_i,l_i,j_i)$ with $\xi_i = \xi_0\,f(n_i,l_i,j_i)$ and
$E_i = 1/\xi_i$, where the absolute values carry the fractal correction
$K_{\mathrm{frak}}$; second, in ladder form
$m_i = r_i\,\xi^{\pi_i}\,v$ with rational $r_i$ and the
generation hierarchy $\pi = \frac32, 1, \frac23$; third, as
$\mathbb{Z}_3$ circulant, $\sqrt{m_k} = M\bigl(1+\sqrt2\cos(\frac29+2\pi k/3)\bigr)$,
with the Koide relation $Q = 2/3$ and the angle $\theta = 2/9$ as
output of the diagonalization. The dimensionless ratios are
exact; deviations arise only in the transition to absolute
SI values, from the SI conversion, the bridge constants $v$ and $E_0$,
the fractal correction, and the finite measurement precision. That is a
structural proof: The ratios follow from the geometry, the
absolute values from the bridge.\status{K}

\section{Critical Appraisal}

\textbf{Ontological status.} The compact time is real and
geometric in the corpus, not merely computational technique. The boundary $\partial T^4 = \emptyset$
makes the question of an initial time point not wrongly answered,
but wrongly posed. Time is discrete: The FFGFT possesses a
minimal, by $\xi$ fixed time increment $d_1 = 100\xi_0 = 1/75$,
and the admissible states are discrete modes. Wheeler's
\enquote{time as illusion} is not shared by the rolled-up view: It
takes time as real, but geometrically different than accustomed.

\textbf{The absolute scale: fixed, not free.} The mass ratios
follow from the mode structure and are exact. The absolute scale is
not a free parameter: The ratio framework rests on $\xi_0 = 4/30000$,
the absolute values carry their SI correction via fixed bridge constants
--- $v$ in the masses, $E_0 = 7{,}398$\,MeV in $\alpha = \xi E_0^2$,
fixed numerical factors in $G$. The Planck quantities are not an input,
but follow from $\xi$: $\xi\to\{m_e,\dots\}\to G\to\ell_P$, and
$\hbar$ itself is geometric. The only exact $\xi$ power relation
is $L_0/\ell_P = \xi_0$, $L_0\approx2{,}15\times10^{-39}$\,m. The
time cycle is not coupled to the Planck scale --- the causal direction
is reversed. No modulus remains open to be calibrated; the
dilaton analogy to string theory does not carry here.

\textbf{Causality.} Does a periodic time generate closed
timelike curves? The carrier is a torus with $\pi_1\neq0$ and stable
windings, not a simple circle; the relevant connection is the
topological connection on $T^4$ (Chapter~\ref{ch:verschraenkung}).
The classical CTC problem of the Gödel solutions concerns the
decompactified Lorentz sector; the compact mass-time direction is
spectral, not causally closed.

\textbf{Experimental verification.} The rolled-up view predicts no
hidden heavy mode tower above the accelerator reach. The discrete
spectrum of compact time \emph{is} already the observed particle spectrum via
$\tilde T\cdot m = 1$; there is no additional Kaluza-Klein tower.
Testable are: the particle masses themselves as $\xi$ derivations
(lepton structure, Koide at $10^{-5}$,
$g-2$); the mode index conservation as coupling selection rule; the
CHSH amplitude of order $\xi$; the fractal dimension $D_f$ and the
fundamental length $L_0$. For a clean falsification,
ratio tests alone are not sufficient --- only a forward-fixed,
independently testable prediction is robust.

\section{The Hidden Inconsistency of Standard QM}

Instantaneity is not a QM-specific problem, but an inheritance of
classical physics. In Newtonian mechanics, gravitation acts
instantaneously; the Coulomb potential is instantaneous; Lagrange and
Hamilton mechanics describe the state $(q(t),p(t))$ at a fixed
$t$ --- the calculation is a timeless algebraic evaluation. No one
asks how long this calculation takes, because calculations in this
language have no duration. Pauli only made explicit what was classically silent.

QM inherits this and adds a tension. \textbf{The state
$\psi(t)$ is instantaneous}: a timeless quantity at a fixed $t$, no
duration, no process. \textbf{The evaluation $|\psi(t)|^2$ takes time}:
To verify a probability density, one needs
many measurements, thus real physical time. The formalism is
instantaneous, its verification is temporally extended. The time that
is needed for the probabilistic evaluation stands outside the
formalism --- it is neither operator nor variable, but what
the experimenter expends to even know the state.

Superposition $c_1|\psi_1\rangle+c_2|\psi_2\rangle$ is thus doubly
timeless: as a state (no time operator, no intrinsic
\enquote{when}) and as a statement (probabilistic, needs time for verification, which lies outside).
What exists as superimposed can only be evidenced through repeated measurements over real time --- and
every measurement destroys the superposition.

\section{Three Consequences That Are Rarely Stated}

In textbooks, time is introduced as a parameter without emphasizing
what the absence of a time operator means. It means concretely:

\textbf{Superposition is timeless.} An electron in the state
$c_\uparrow|\uparrow\rangle+c_\downarrow|\downarrow\rangle$ has no
intrinsic time scale. The question \enquote{how long is it in this
superposition?} is not poseable in the formalism. What exists is
the coherence time $T_2$ --- a phenomenological quantity of the
environment coupling, not an intrinsic property of the state. In the
FFGFT, each mode carries $\tilde T_i = 1/m_i$ as a geometric quantity. A
state \emph{is} its time scale --- not as the duration of a
superposition, but as a structural property of the mode.

\textbf{Entanglement without simultaneity.} Two entangled
particles share a state in $\mathcal{H}_1\otimes\mathcal{H}_2$
without built-in time structure. The question of the simultaneity
of a measurement on both is not a QM question, but a classical
one about the external parameter --- and thus a relativistic one. That is
the deepest reason why QM is compatible with Lorentz invariance:
Simultaneity stands outside the formalism. In the FFGFT: If
two modes with $m_1\neq m_2$ were entangled, they would have different
$\tilde T_1\neq\tilde T_2$. Cross-mode entanglement generates
incompatible time scales --- structurally more problematic than in QM,
where the conflict does not arise because $t$ stands outside.

\textbf{Collapse without duration.} The transition $|\psi\rangle\to|\psi_k\rangle$
has no duration in the formalism; it is instantaneous in $t$. How long the
collapse takes lies outside standard QM. Without a time operator there
is no observable \enquote{collapse duration}. In the FFGFT, the
Type I step (measurement act) is a geometric transition: The mass circle
$S^1_m$ becomes the time axis $\mathbb{R}_t$, the winding number is
discarded. This is an information loss with a natural direction.

\section{What $\tilde T\cdot m = 1$ Structurally Changes}

The fundamental relation sets time not as an external parameter, but as a
geometric quantity of each mode: The period $R_m = \hbar/(mc^2)$ is
intrinsic --- it belongs to the particle like its mass. Three
consequences:

\textbf{Superposition} within a mode --- intermediate torsion
--- has a natural time scale $\tilde T_i$. Superposition
\emph{over} different modes would superimpose incompatible time scales;
that is the structural reason why cross-mode superpositions are problematic in the FFGFT.\status{B}

\textbf{Entanglement} within a mode has simultaneity defined by
$\tilde T$. Across modes, a common time concept is lacking ---
similar to QM, but for a different reason: not
because $t$ stands externally, but because different modes carry different
geometric time scales.

\textbf{Collapse} (Type I) has a natural direction ---
lossy from $S^1_m$ to $\mathbb{R}_t$ --- and a natural
asymmetry: The reverse direction requires the discarded winding number. The
time arrow is not an additional assumption, but a consequence of the geometry of
measurement.

The Pauli theorem applies to an open time parameter. For a
compact time dimension it does not apply: Its conjugate has a
discrete spectrum, the winding numbers. $\tilde T\cdot m = 1$
realizes exactly this situation: Time is compact, the conjugate
spectrum is the discrete mass ladder. The theorem does not exclude this case
--- it enforces it as the only consistent way to integrate time into
the Hilbert space.

\section{Unrolling: Type I and Type III}

Two different operations carry the same name. \textbf{Type I
unrolling} is the measurement act: lossy, the projection from the
compact carrier onto the real axis is not reversible. That is the
collapse. \textbf{Type III unrolling} is the representation change: the
bijective pullback from Chapter~\ref{ch:hilbert}, no information loss,
only a language change. The correction factor $K_{\mathrm{frak}}$ sits in the
mass formula of the Type I transition; it is not involved in the Type III pullback. Anyone who calculates in Hilbert space language sees $K_{\mathrm{frak}}$
only in the SI translation.\status{B}

\section{Summary}

\begin{center}
\small\renewcommand{\arraystretch}{1.3}
\begin{tabular}{p{2.2cm}p{3.5cm}p{3.7cm}}
\toprule
 & \textbf{Standard QM} & \textbf{FFGFT} \\
\midrule
Time & absent: no operator; external parameter & geometric: $\tilde T_i = 1/m_i$, compact, mode-bound \\
Superposition & timeless, no intrinsic clock & intermediate torsion in a mode; $\tilde T_i$ as scale \\
Cross-mode & algebraically allowed, no time conflict & incompatible $\tilde T_i\neq\tilde T_j$, structurally problematic \\
Entanglement & timeless; simultaneity external (Lorentz) & mode-specific $\tilde T$; no common time concept across modes \\
Collapse & instantaneous in $t$, no duration & Type I: directed information loss, time arrow geometric \\
Pauli theorem & blocks time operator for open $t$ & does not apply for compact $\tilde T$ \\
\bottomrule
\end{tabular}
\end{center}

The choice of where time sits precedes the mathematics and decides
whether the modes are given or attached. Anyone who reads the sectors as
spatial eigenspaces keeps the familiar physics and pays with
the external attachment of the numbers. Anyone who takes time as a closed
dimension into the space obtains the modes structurally and pays with
an ontological commitment. The ontological price of the second choice
remains on the table --- but the methodological price of the first is not
smaller, it is only more familiar.

\quelle{Doc.~307 (Time in state space, July 2026), Doc.~334 (Superposition without time, Aug.\ 2026), Doc.~333 (Type I/III), Doc.~306 ($T\cdot E = 1$), Doc.~270 (Decompactification)}